\section{Infrastructure Design}
\label{sec:infra}

Block AttnRes introduces additional system challenges compared to standard residual connections. For large-scale model training, block representations must be propagated across pipeline stages, causing heavy communication in a na\"ive implementation. During inference, repeated access to accumulated block representations increases latency, while long-context prefilling amplifies the memory cost of caching block representations. We address these challenges with cross-stage caching in training, and with a two-phase computation strategy together with a memory-efficient prefilling scheme in inference.

\subsection{Training}
\label{sec:infra-training}
For small-scale training, AttnRes adds a tiny computation overhead and no extra memory usage, as the activations need to be saved for backpropagation regardless.
Under large-scale distributed training, pipeline parallelism poses the primary infrastructure challenge for AttnRes.
Full AttnRes requires all $L$ layer outputs to be transmitted across stages; Block AttnRes reduces this to $N$ block representations, and the optimizations below further minimize the remaining overhead.

\paragraph{Pipeline communication.}
With standard residual connections, pipeline parallelism~\citep{huang2019gpipe} transfers a fixed-size hidden state between adjacent stages, independent of pipeline depth.
Block AttnRes requires all accumulated block representations at each stage for inter-block attention, and na\"ively transmitting the full history at every transition incurs redundant communication.

Consider an interleaved pipeline schedule~\citep{narayanan2021megatron} with $P$ physical stages and $V$ virtual stages per physical stage.
For simplicity, assume each physical stage produces on average $N_p$ block representations of dimension $d$ per token.\footnote{In practice, block boundaries need not align with physical stage boundaries. For example, in Fig.~\ref{fig:pipeline}, each block spans two physical stages, so only every other transition involves a newly completed block.}
With $C = PV$ total chunks (each physical stage in each virtual stage), the $j$-th chunk accumulates $jN_p$ blocks.
Na\"ively transmitting all accumulated blocks at every transition incurs per-token communication cost:
\begin{equation}
    \mathrm{Comm}_{\text{na\"ive}} = \sum_{j=1}^{C-1} jN_p \cdot d = \frac{C(C{-}1)}{2}\,N_p d.
\end{equation}
\paragraph{Cross-stage caching.}
Since each physical stage processes multiple virtual stages in succession, we can eliminate this redundancy by caching blocks locally: blocks received during earlier virtual stages remain in local memory and need not be re-transmitted.
The first virtual stage ($v = 1$) has no cache and accumulates normally; for $v \geq 2$, each transition conveys only the ${\sim}P N_p$ incremental blocks accumulated since the receiver's corresponding chunk in the previous virtual stage.
Total communication reduces to:
\begin{equation}
    \mathrm{Comm}_{\text{cached}} = \underbrace{\frac{P(P{-}1)}{2}\, N_p d}_{\text{first virtual stage}} + \underbrace{(V{-}1)\, P^2\, N_p d}_{\text{subsequent virtual stages}}.
\end{equation}

\begin{figure*}[t]
    \centering
    \pgfdeclarelayer{arrowlayer}
    \pgfdeclarelayer{boxlayer}
    \pgfsetlayers{background,arrowlayer,main,boxlayer}
    \begin{adjustbox}{width=0.99\textwidth,center}
        \begin{tikzpicture}[
                box/.style={
                        rectangle,
                        draw,
                        very thick,
                        rounded corners=2pt,
                        minimum width=0.8cm,
                        minimum height=0.8cm,
                        font=\normalsize
                    },
                round0/.style={box, fill=red!20, draw=black},
                round1/.style={box, fill=blue!15, draw=black},
                round3/.style={box, fill=teal!35, draw=black},
                round0pattern/.style={box, preaction={fill=white}, fill=red!20, draw=black, pattern=north east lines, pattern color=red!60},
                round1pattern/.style={box, preaction={fill=white}, fill=blue!15, draw=black, pattern=north east lines, pattern color=blue!50},
                round3pattern/.style={box, preaction={fill=white}, fill=teal!35, draw=black, pattern=north east lines, pattern color=teal!70},
                label/.style={above},
                cache/.style={below},
                rank/.style={anchor=east},
                arrow/.style={->, >=stealth, line width=1pt, rounded corners=4pt, shorten >=3pt, shorten <=3pt}
            ]

            \node[rank] at (-0.5, 0) {\textsc{Rank 0}};
            \node[rank] at (-0.5, -1.5) {\textsc{Rank 1}};
            \node[rank] at (-0.5, -3) {\textsc{Rank 2}};
            \node[rank] at (-0.5, -4.5) {\textsc{Rank 3}};


            \node[fill=white, text width=60pt, inner sep=10pt, align=left] (cache0r1) at (2.0, 0) {$[\;\bm{b}_0\;]$};
            \node[fill=white, text width=25pt, inner sep=15pt, align=left] (cache0br1) at ([xshift=0pt]cache0r1.east) {$[\;\phantom{\bm{b}_0}\;]$};
            \node[fill=white, text width=60pt, inner sep=10pt, align=left] (cache0r2) at (2.0, -1.5) {$[\;\bm{b}_0\;]$};
            \node[fill=white, text width=25pt, inner sep=15pt, align=left] (cache0br2) at ([xshift=0pt]cache0r2.east) {$[\;\bm{b}_1\;]$};
            \node[fill=white, text width=60pt, inner sep=10pt, align=left] (cache0r3) at (2.0, -3) {$[\;\bm{b}_0,\bm{b}_1\;]$};
            \node[fill=white, text width=25pt, inner sep=15pt, align=left] (cache0br3) at ([xshift=0pt]cache0r3.east) {$[\;\phantom{\bm{b}_0}\;]$};
            \node[fill=white, text width=60pt, inner sep=10pt, align=left] (cache0r4) at (2.0, -4.5) {$[\;\bm{b}_0,\bm{b}_1\;]$};
            \node[fill=white, text width=25pt, inner sep=15pt, align=left] (cache0br4) at ([xshift=0pt]cache0r4.east) {$[\;\bm{b}_2\;]$};

            \begin{pgfonlayer}{boxlayer}
                \node[round0] (r0c0) at ([xshift=35pt]cache0br1.east) {1};
                \fill[white] ([xshift=63pt]cache0br1.east) ++(-0.4cm,-0.4cm) rectangle ++(0.8cm,0.8cm);
                \node[round0, opacity=0.3] (r0c0b) at ([xshift=63pt]cache0br1.east) {2};
                \node[round0pattern] (r1c0) at ([xshift=63pt]cache0br2.east) {1};
                \fill[white] ([xshift=91pt]cache0br2.east) ++(-0.4cm,-0.4cm) rectangle ++(0.8cm,0.8cm);
                \node[round0pattern, opacity=0.3] (r1c0b) at ([xshift=91pt]cache0br2.east) {2};
                \node[round0] (r2c0) at ([xshift=91pt]cache0br3.east) {1};
                \fill[white] ([xshift=119pt]cache0br3.east) ++(-0.4cm,-0.4cm) rectangle ++(0.8cm,0.8cm);
                \node[round0, opacity=0.3] (r2c0b) at ([xshift=119pt]cache0br3.east) {2};
                \node[round0pattern] (r3c0) at ([xshift=119pt]cache0br4.east) {1};
                \fill[white] ([xshift=147pt]cache0br4.east) ++(-0.4cm,-0.4cm) rectangle ++(0.8cm,0.8cm);
                \node[round0pattern, opacity=0.3] (r3c0b) at ([xshift=147pt]cache0br4.east) {2};
            \end{pgfonlayer}

            \node[fill=white, text width=60pt, inner sep=10pt, align=left] (cache1r1) at ([xshift=120pt]r3c0b.east |- cache0r1) {$+\; [\;\bm{b}_1, \bm{b}_2\;]$};
            \node[fill=white, text width=25pt, inner sep=15pt, align=left] (cache1br1) at ([xshift=0pt]cache1r1.east) {$[\;\phantom{\bm{b}_3}\;]$};
            \node[fill=white, text width=60pt, inner sep=10pt, align=left] (cache1r2) at ([xshift=120pt]r3c0b.east |- cache0r2) {$+\; [\;\bm{b}_1, \bm{b}_2\;]$};
            \node[fill=white, text width=25pt, inner sep=15pt, align=left] (cache1br2) at ([xshift=0pt]cache1r2.east) {$[\;\bm{b}_3\;]$};
            \node[fill=white, text width=60pt, inner sep=10pt, align=left] (cache1r3) at ([xshift=120pt]r3c0b.east |- cache0r3) {$+\; [\;\bm{b}_2,\bm{b}_3\;]$};
            \node[fill=white, text width=25pt, inner sep=15pt, align=left] (cache1br3) at ([xshift=0pt]cache1r3.east) {$[\;\phantom{\bm{b}_4}\;]$};
            \node[fill=white, text width=60pt, inner sep=10pt, align=left] (cache1r4) at ([xshift=120pt]r3c0b.east |- cache0r4) {$+\; [\;\bm{b}_2,\bm{b}_3\;]$};
            \node[fill=white, text width=25pt, inner sep=15pt, align=left] (cache1br4) at ([xshift=0pt]cache1r4.east) {$[\;\bm{b}_4\;]$};

            \begin{pgfonlayer}{boxlayer}
                \node[round1] (r0c1) at ([xshift=35pt]cache1br1.east) {1};
                \fill[white] ([xshift=63pt]cache1br1.east) ++(-0.4cm,-0.4cm) rectangle ++(0.8cm,0.8cm);
                \node[round1, opacity=0.3] (r0c1b) at ([xshift=63pt]cache1br1.east) {2};
                \node[round1pattern] (r1c1) at ([xshift=63pt]cache1br2.east) {1};
                \fill[white] ([xshift=91pt]cache1br2.east) ++(-0.4cm,-0.4cm) rectangle ++(0.8cm,0.8cm);
                \node[round1pattern, opacity=0.3] (r1c1b) at ([xshift=91pt]cache1br2.east) {2};
                \node[round1] (r2c1) at ([xshift=91pt]cache1br3.east) {1};
                \fill[white] ([xshift=119pt]cache1br3.east) ++(-0.4cm,-0.4cm) rectangle ++(0.8cm,0.8cm);
                \node[round1, opacity=0.3] (r2c1b) at ([xshift=119pt]cache1br3.east) {2};
                \node[round1pattern] (r3c1) at ([xshift=119pt]cache1br4.east) {1};
                \fill[white] ([xshift=147pt]cache1br4.east) ++(-0.4cm,-0.4cm) rectangle ++(0.8cm,0.8cm);
                \node[round1pattern, opacity=0.3] (r3c1b) at ([xshift=147pt]cache1br4.east) {2};
            \end{pgfonlayer}

            \begin{scope}[on background layer]
                \foreach \y in {0, -1.5, -3, -4.5} {
                        \draw[dotted, gray] ([xshift=-10pt]cache0r1.west |- 0,\y) -- ([xshift=10pt]r3c1b.east |- 0,\y);
                    }
            \end{scope}

            \coordinate (round1y) at ([yshift=15pt]cache0r1.north);
            \draw[line width=2pt] (cache0r1.north west |- round1y) -- ([xshift=10pt]r3c0b.north east |- round1y) node[midway, above=3pt] {\textsc{Virtual Stage 0}};
            \coordinate (round2y) at ([yshift=15pt]cache1r1.north);
            \draw[line width=2pt] (cache1r1.north west |- round2y) -- ([xshift=10pt]r3c1b.north east |- round2y) node[midway, above=3pt] {\textsc{Virtual Stage 1}};

            \begin{pgfonlayer}{arrowlayer}
                \draw[arrow] (r0c0.west) -- ++(-0.8,0) |- ([yshift=0pt]r1c0.west);
                \draw[arrow] ([yshift=-10pt]r1c0.west) -- ++(-0.8,0) |- ([yshift=0pt]r2c0.west);
                \draw[arrow] ([yshift=-10pt]r2c0.west) -- ++(-0.8,0) |- (r3c0.west);
                \draw[arrow] (r3c0.east) -- ($(r3c0b.east)!0.5!(cache1r1.west |- r3c0.east)$) |- (r0c1.west);
                \draw[arrow] ([yshift=-10pt]r0c1.west) -- ++(-0.8,0) |-     ([yshift=0pt]r1c1.west);
                \draw[arrow] ([yshift=-10pt]r1c1.west) -- ++(-0.8,0) |- ([yshift=0pt]r2c1.west);
                \draw[arrow] ([yshift=-10pt]r2c1.west) -- ++(-0.8,0) |- (r3c1.west);
            \end{pgfonlayer}

        \end{tikzpicture}
    \end{adjustbox}
    \caption{
    Cache-based pipeline communication example with 4 physical ranks and 2 virtual stages per rank, where hatched boxes denote end of AttnRes blocks.
    Numbers indicate micro-batch indices. Each rank caches previously received blocks; stage transitions only transmit incremental blocks ($+[\bm{b}_1, \bm{b}_2]$) instead of the full history.
    }
    \label{fig:pipeline}
\end{figure*}


Caching reduces peak per-transition cost from $O(C)$ to $O(P)$, a $V{\times}$ improvement that enables full overlap with computation during steady-state 1F1B.
The backward pass benefits from the same scheme. Fig.~\ref{fig:pipeline} illustrates this optimization with $P{=}4$ and $V{=}2$: for the second virtual stage, caching eliminates 6 redundant block transmissions.

\paragraph{Memory overhead.}
With cross-stage caching, each block is stored exactly once across all $V$ virtual stages, which becomes negligible relative to standard per-layer activation cache. Crucially, the per-layer activation footprint remains identical to standard architectures, as activation checkpointing eliminates all inter-block attention intermediates, and the checkpointed input $\bm{p}_l$ matches the memory size of the hidden state $\bm{h}_l$ it replaces.

In terms of wall-clock time, Block AttnRes adds negligible training overhead when pipeline parallelism is not enabled; under pipeline parallelism, the measured end-to-end overhead is less than 4\%.





















\subsection{Inference}
\label{sec:infra-inference}
The two-phase computation strategy described below applies to both Full and Block AttnRes: in either case, layers are grouped into blocks of size $S$, with Phase~1 batching the inter-block queries and Phase~2 handling sequential intra-block lookback. For Full AttnRes, this reduces per-layer I/O from $O(Ld)$ to $O((S{+}N)d)$ (detailed derivation shown in Appendix \ref{app:inference}); Block AttnRes further reduces the stored representations from $L$ to $N$, since each block is compressed into a single vector.
In what follows, we focus on Block AttnRes and detail the two-phase computation strategy together with a sequence-sharded prefilling scheme for long-context inputs.

\paragraph{Two-phase computation strategy.}
The layer-wise attention computation of Block AttnRes resembles autoregressive decoding, where block representations serve as a shared KV cache reused across layers.
A na\"ive implementation computes the attention residual at every layer, each requiring a full pass over all preceding blocks, resulting in $O(L \cdot N)$ memory accesses.
Since the pseudo-query vectors are decoupled from the forward computation (\S\ref{sec:attnres}), all $S = L/N$ queries within a block can be batched into a single matrix multiplication, amortizing memory access from $S$ reads to $1$.

Algorithm~\ref{alg:inference} instantiates a two-phase computation strategy exploiting this property.
\begin{itemize}[leftmargin=*]
    \item \textbf{Phase~1} computes inter-block attention for all $S$ layers simultaneously via a single batched query against the cached block representations, returning both outputs and $\operatorname{softmax}$ statistics (max and log-sum-exp). This amortizes the memory access cost, reducing reads from $S$ times to just once per block.
    \item \textbf{Phase~2} computes intra-block attention sequentially for each layer using the evolving partial sum, then merges with Phase~1 outputs through online $\operatorname{softmax}$~\citep{milakov2018online}. Because the online-$\operatorname{softmax}$ merge is elementwise, this phase naturally admits kernel fusion with surrounding operations, further reducing I/O overhead.
\end{itemize}

\begin{algorithm}[t]
    \caption{Two-phase computation for block $n$}
    \label{alg:inference}
    \KwIn{Pseudo queries $\{\bm{w}_l\}_{l \in \mathcal{B}_n}$, block representations $\{\bm{b}_0, \ldots, \bm{b}_{n-1}\}$}
    \BlankLine
    \tcc{Phase 1: Parallel inter-block attention}
    $\mathbf{Q} \gets [\bm{w}_l]_{l \in \mathcal{B}_n}$\tcp*{$[S, d]$}
    $\mathbf{K}, \mathbf{V} \gets [\bm{b}_0; \ldots; \bm{b}_{n-1}]$\tcp*{$[n, d]$}
    $\{\bm{o}_l^{(1)}, m_l^{(1)}, \ell_l^{(1)}\}_{l \in \mathcal{B}_n} \gets \textsc{AttnWithStats}(\mathbf{Q}, \mathbf{K}, \mathbf{V})$\tcp*{Return LSE}
    \;
    \tcc{Phase 2: Sequential intra-block attention + Online $\operatorname{softmax}$ merge}
    $i \gets 0$\;
    \For{$l \in \mathcal{B}_n$}{
        \eIf{$i = 0$}{
            $\bm{h}_l \gets \bm{o}_l^{(1)} / \ell_l^{(1)}$\tcp*{Inter-block only}
        }{
            $\bm{o}_l^{(2)}, m_l^{(2)}, \ell_l^{(2)} \gets \textsc{AttnWithStats}(\bm{w}_l, \bm{b}_n^i, \bm{b}_n^i)$\tcp*{Intra-block}
            $m_l \gets \max(m_l^{(1)}, m_l^{(2)})$\;
            $\bm{h}_l \gets \dfrac{e^{m_l^{(1)} - m_l} \bm{o}_l^{(1)} + e^{m_l^{(2)} - m_l} \bm{o}_l^{(2)}}{e^{m_l^{(1)} - m_l} \ell_l^{(1)} + e^{m_l^{(2)} - m_l} \ell_l^{(2)}}$\tcp*{Online softmax merge}
        }
        $i \gets i + 1$\;
        $\bm{b}_n^{i} \gets \bm{b}_n^{i-1} + f_l(\bm{h}_l)$\tcp*{Update partial sum; $\bm{b}_n^0 \coloneqq \bm{0}$}
    }
    \KwRet{$\{\bm{h}_l\}_{l \in \mathcal{B}_n}$}
\end{algorithm}

With the two-phase design, Phase~2 preserves an I/O footprint similar to that of standard residual connections, whereas the main additional cost arises from Phase~1 inter-block attention. Because these inter-block reads are amortized across all layers in a block through batching, the total per-layer memory access cost remains only $(\frac{N}{S} + 3)d$ reads and $2d$ writes (Table~\ref{tab:memory_access}). This is substantially lower than the residual-stream I/O of prior residual generalizations such as (m)HC under typical settings. In practice, Phase~1 can also partially overlap with the computation of the first layer in the block, further reducing its wall-clock impact. As a result, the end-to-end inference latency overhead is less than 2\% on typical inference workloads.

\begin{table}[h]
  \centering
  \small
  \setlength{\tabcolsep}{9pt}
  \renewcommand{\arraystretch}{1.1}
  \caption{Memory access cost per token per layer incurred by the residual mechanism under each scheme. The internal I/O of the layer function $f_l$ is excluded. For AttnRes, both Full and Block variants use the two-phase inference schedule described in Appendix~\ref{app:inference}; amortized costs are averaged over $N$ layers within a block. Typical values: $L{=}128$, $N{=}8$, $S{=}L/N{=}16$, $m{=}4$.}
  \label{tab:memory_access}
  \newlength{\iocolwidth}\settowidth{\iocolwidth}{$(8m{+}2)d{+}2m^2{+}4m$}%
  \begin{tabular}{@{}l l l c c w{c}{\iocolwidth} w{c}{\iocolwidth}@{}}
    \toprule
                                                           &                                                               & \multirow{2}{*}{Operation} & \multirow{2}{*}{Read}           & \multirow{2}{*}{Write}                   & \multicolumn{2}{c}{Total I/O}                                                \\
    \cmidrule(l){6-7}
                                                           &                                                               &                            &                                 &                                          & Symbolic                                           & Typical                 \\
    \midrule
    \multicolumn{2}{l}{Standard Residuals}                 & Residual Merge                                                & $2d$                       & $d$                             & $3d$                                     & $3d$                                                                         \\
    \midrule
    \multicolumn{2}{l}{\multirow{5}{*}{mHC ($m$ streams)}} & Compute $\bm{\alpha}_{l}$, $\bm{\beta}_{l}$, $\mathbf{A}_{l}$ & $md$                       & $m^2{+}2m$                      & \multirow{5}{*}{$(8m{+}2)d{+}2m^2{+}4m$} & \multirow{5}{*}{$34d$}                                                       \\
    \multicolumn{2}{l}{}                                   & Apply $\bm{\alpha}_{l}$                                       & $md{+}m$                   & $d$                             &                                          &                                                                              \\
    \multicolumn{2}{l}{}                                   & Apply $\bm{\beta}_{l}$                                        & $d{+}m$                    & $md$                            &                                          &                                                                              \\
    \multicolumn{2}{l}{}                                   & Apply $\mathbf{A}_{l}$                                        & $md{+}m^2$                 & $md$                            &                                          &                                                                              \\
    \multicolumn{2}{l}{}                                   & Residual Merge                                                & $2md$                      & $md$                            &                                          &                                                                              \\
    \midrule
    \multirow{4}{*}{AttnRes}                               %& Full (na\"ive)                                                & Attention (amortized)             & $(L{+}1)d$                          & $d$                                      & $(L{+}2)d$                                         & $130d$                  \\
                                                           & \multirow{2}{*}{Full}                                         & Phase~1 (amortized)        & $(N{-}1)d$                      & $d$                                      & \multirow{2}{*}{$(S{+}N)d$}                        & \multirow{2}{*}{$24d$}  \\
                                                           &                                                               & Phase~2                    & $(S{-}1)d$                      & $d$                                      &                                                    &                         \\
    \cmidrule{2-7}
                                                           & \multirow{2}{*}{Block}                                        & Phase~1 (amortized)        & $\frac{N}{S}d$                  & $d$                                      & \multirow{2}{*}{$\left(\frac{N}{S}{+}5\right)d$}   & \multirow{2}{*}{$5.5d$} \\
                                                           &                                                               & Phase~2                    & $3d$                            & $d$                                      &                                                    &                         \\
    \bottomrule
  \end{tabular}
\end{table}




\paragraph{Memory-efficient prefilling.}
Storing block representations during prefilling requires $N \cdot T \cdot d$ elements, which incurs 15\,GB of memory for a 128K-token sequence with 8 blocks. We mitigate this by sharding these representations along the sequence dimension across $P$ tensor-parallel devices, allowing Phase~1 to execute independently on local sequence shards. The Phase~2 online-$\operatorname{softmax}$ merge then integrates into the standard TP all-reduce communication path: the output is reduce-scattered, merged locally, and reconstructed via all-gather, naturally admitting kernel fusion with operations like $\operatorname{RMSNorm}$. This reduces the per-device memory footprint to $N \cdot (T/P) \cdot d$---lowering the 128K-context example from 15\,GB to roughly 1.9\,GB per device. Combined with chunked prefill (e.g., 16K chunk size), the overhead further reduces to under 0.3\,GB per device.


